% 可选论文组合图：放在main.tex的\end{document}之前。中文编译器与原工程一致。

\begin{figure}[p]\centering
\includegraphics[width=.98\linewidth,height=.80\textheight,keepaspectratio]{figures/P01.pdf}
\caption{地形—载荷—能耗的条件响应。上图与两个截面来自同一确定性模型网格。固定S008距离和端点，仅基准巡航海拔为题定任务；不是现场数据或新调度结果。}
\label{fig:P01}
\end{figure}

\begin{figure}[p]\centering
\includegraphics[width=.98\linewidth,height=.80\textheight,keepaspectratio]{figures/P02.pdf}
\caption{返航安全余量与运输能力。小提琴为每档15个服务区的C型确定性横截面，实点与四分位线均来自结果。敏感性从20\%开始；不可行档位不记为0。}
\label{fig:P02}
\end{figure}

\begin{figure}[p]\centering
\includegraphics[width=.98\linewidth,height=.80\textheight,keepaspectratio]{figures/P03.pdf}
\caption{需求空间分布与物资组成。地图仅呈现真实离散服务节点，不制造连续需求热区；矩阵由80个真实货箱汇总。累计需求按期望时刻统计，不是实际交付或到达率。}
\label{fig:P03}
\end{figure}

\begin{figure}[p]\centering
\includegraphics[width=.98\linewidth,height=.80\textheight,keepaspectratio]{figures/P04.pdf}
\caption{已计算方案的多指标比较。上图只含5个真实Q3候选，其中有同点重合；下图为已存的候选和搜索轮次，不扩充为虚假的种群点云或全局帕累托前沿。}
\label{fig:P04}
\end{figure}

\begin{figure}[p]\centering
\includegraphics[width=.98\linewidth,height=.80\textheight,keepaspectratio]{figures/P05.pdf}
\caption{单架次的时空—能源—通信状态。4幅图共享同一绝对时间范围，来源为Q3-T-016。SOC按模型分阶段累计并校验返航端点；末段显示电池充满，不把充电结束当作Q3完工。}
\label{fig:P05}
\end{figure}

\begin{figure}[p]\centering
\includegraphics[width=.98\linewidth,height=.80\textheight,keepaspectratio]{figures/P06.pdf}
\caption{严格三分区与独立资源需求。保持最终Q3任务和通信关系不变。三组需求35件、分类型缺口7件；资源总件数不是采购成本，补足型号资源才可执行。}
\label{fig:P06}
\end{figure}
